\section{14. Death and resurrection}

Character death is meaningful, but a powerful recovery mechanism prevents early deaths from permanently excluding players.

\subsection{Resurrection Runes}

A Resurrection Rune is awarded when a boss is defeated and may appear in chests after all unique challenge rewards are exhausted. When activated, it resurrects all other dead player characters at full health and returns their items. It is consumed on use. If its bearer dies, it automatically transfers to another living player character. Living characters elsewhere are not transported or otherwise affected. Resurrected characters can be placed in the activating character's chamber using the respawn functionality.

\subsection{Recommended adjudication}

Treat the rune as powerful magical recovery rather than a conventional resurrection spell. Full health need not restore expended spell slots, class abilities or consumed magical items. Avoid duplicating equipment already recovered by companions. Determine activation action and timing consistently; requiring an action during combat is reasonable.

\subsection{Total party death}

If no living character remains, nobody can activate a rune. Unless the DM deliberately introduced another recovery mechanism, this constitutes defeat. Do not automatically rescue the party simply because the adventure is dangerous. The possibility of failure gives decisions meaning.
